\section{问题的提出：需求数据能告诉我们什么}
% ============================================================
% 第六讲 6.1：反问题的三种提法与需求函数的信息含量
% ============================================================

\subsection{正向链条与它的三个反问题}

前五章建立的是一条单向的推导链。用一张图概括：

\begin{center}
\begin{tikzpicture}[>=Stealth]
  \node[boxsoft] (pref) at (0,0) {偏好 $\succsim$};
  \node[boxsoft] (util) at (4.1,0) {效用 $u(\cdot)$};
  \node[boxsoft] (dem)  at (8.5,0) {需求 $x(p,w)$};
  \node[boxsoft] (dual) at (8.5,-2.7) {支出函数 $e(p,\bar u)$、\\[1pt] Hicksian 需求 $h(p,\bar u)$};
  \draw[arrblue] (pref)--node[figlabel,above,align=center]{表示定理\\[-2pt]{\scriptsize Debreu}} (util);
  \draw[arrblue] (util)--node[figlabel,above,align=center]{最优化\\[-2pt]{\scriptsize Kuhn--Tucker}} (dem);
  \draw[arrblue] (dem)--node[figlabel,right,align=left]{对偶\\[-2pt]{\scriptsize Shephard 引理}} (dual);
  \node[figlabel,below=0.18cm of dual,align=center]
       {Slutsky 矩阵：$S=D_p^2e=\nabla_ph$};
  \draw[thick,dashed,->] (dem.north) to[bend right=26]
       node[figlabel,above,align=center]{反问题：可积性（本章）} (pref.north);
\end{tikzpicture}
\end{center}
\figexplan{实线为前五章建立的正向链条：偏好经表示定理得到效用，经最优化得到需求，再经对偶得到支出函数与 Hicksian 需求（支出函数的 Hessian 即 Slutsky 矩阵，这是全章的枢纽）。上方那条虚线 arc 是本章要回答的反问题：从最右端的 $x(p,w)$ 能否退回最左端的偏好，退回时丢失了什么。注意反问题不必经过效用函数——\S 6.4 的路线是直接把需求积分成支出函数，再由对偶反造偏好。}

反问题可以精确化为三个层次，难度依次上升：

\begin{enumerate}
\item \textbf{函数层面}（可积性）。已知处处光滑的需求函数 $x:\R^n_{++}\times\R_{++}\to\R^n_+$ 的解析式，
      是否存在一个偏好 $\succsim$ 使 $x(p,w)$ 恰为它在预算约束下的需求？
\item \textbf{数据层面}（可理性化）。只观察到有限个价格—数量对 $\{(p^t,x^t)\}_{t=1}^{T}$，
      是否存在一个效用函数与之一致？注意这里连"需求函数"都不知道，只有点。
\item \textbf{福利层面}（货币度量）。若前两层通过，能否用已经恢复出来的 $e(p,\bar u)$
      计算价格变化对消费者的货币影响（补偿变差、等价变差）？
\end{enumerate}

三个层次分别对应本章的 \S 6.2--6.5、\S 6.6、\S 6.8。三个层次的共同前提是一个方法论立场：
\textbf{理论对数据施加了可检验的限制}。如果任何光滑函数都能充当需求函数，需求估计就不存在"设定检验"这一说；
恰恰因为存在限制，计量经济学家才能用数据去拒绝或接受理论。本章的技术内容，本质上就是把这组限制写清楚、
证明它充分、并说明它在什么维数下才有内容。

\subsection{三条基本观察}

在动笔之前必须先弄清楚：反问题\emph{不可能}有唯一的解。原因有三，逐条陈述。

\notehead{观察一：从偏好到需求是多对一的，因此只能恢复到序数水平。}

设 $f:\R\to\R$ 严格递增，则 $f\circ u$ 与 $u$ 刻画同一偏好，从而生成同一个需求函数。
更一般地，任何保持序关系的单调变换都不改变任何可观测的选择。因此反问题的解集若含一个效用函数，
必含它的一整条单调变换轨道。这不是技术缺陷，而是效用理论的既有性质：效用是序数概念，
需求数据也不可能提供比序数更多的信息。\textbf{能期望的最强结论是"恢复到一个单调变换为止"。}

\begin{example}[CD 的两条等价轨道]
取 $u(x)=x_1^{1/2}x_2^{1/2}$ 与 $u'(x)=x_1x_2$。两者生成同一个需求
\[
 x(p,w)=\Bigl(\frac{w}{2p_1},\;\frac{w}{2p_2}\Bigr),
\]
因为 $u'=(u)^2$ 是单调变换。但若以 $u$ 为效用的度量，$e(p,\bar u)=\bar u\,(4p_1p_2)^{1/2}$；
以 $u'$ 为度量，$e'(p,\bar u')=\bar u'\,(p_1p_2)^{1/2}$。两个支出函数不同，然而它们描述的是同一组无差异曲线。
这提示：\textbf{支出函数的值依赖效用标度，但支出函数的两阶导数（即 Slutsky 矩阵）不依赖}——
后者才是需求数据能确定的东西。
\end{example}

\notehead{观察二："序数"不等于"没有内容"，可确定的是曲率而非水平。}

虽然效用的绝对水平不可识别，但无差异曲线的\emph{形状}是唯一的（在同一偏好类内）。
形状在支出函数上的体现是它对价格的二阶导数：$D_p^2e=\nabla_p h$。
这正是可观测的 Slutsky 矩阵。因此本讲的核心结论可以先用一句话预告：
\textbf{需求数据唯一确定偏好的序数结构，其可检验内容集中在 Slutsky 矩阵的对称性与负半定性上。}
而且一旦选定效用标度（例如规定 $e(p,\bar u)=\bar u$ 对某个基准价格向量成立），
支出函数就被完全钉死，福利度量随之确定。

\notehead{观察三：需求函数不能是任意的，否则理论不可检验。}

第三条观察最好用一个具体的失败例子说明。考虑两期观测，两商品，价格向量相同：

\begin{example}[同一价格下选两个不同束]\label{ex:twobundles}
设 $p^1=p^2=(1,1)$，观测到 $x^1=(1,0)$、$x^2=(0,1)$。
若这两期数据来自同一个按预算约束最大化的消费者，则两期的财富必须相同（否则预算集不同，价格相同而"同一价格下两个束都可负担"这句话就不成立）：
$w^1=p^1\cdot x^1=1$，$w^2=p^2\cdot x^2=1$。于是在第 1 期，$x^2$ 也恰好可负担（$p^1\cdot x^2=1\le w^1$），
而消费者选了 $x^1$，说明 $u(x^1)\ge u(x^2)$；对称地，第 2 期说明 $u(x^2)\ge u(x^1)$。
两式相容，并不矛盾——但若把两式都取严格（例如观测到 $x^1$ 时 $x^2$ 严格便宜），矛盾立刻出现。
这就是弱显示偏好公理（WARP）的原型，见 \S 6.6。
\end{example}

这个例子还说明一件事：\textbf{理论限制可以被数据违反}，因此可以被检验。
第 5 章已经证明"任何由偏好生成的需求都满足 Slutsky 对称与负半定"；
本章要做的是反过来证明"满足这两条的需求都能由某个偏好生成"（充分性），
从而把这两条确立为\emph{充分必要}的可理性化条件。定理的两半合起来才构成完整的可检验性陈述：
一半告诉我们理论禁止什么，另一半告诉我们理论允许什么。

\subsection{两条恢复路径与一个警告}

从技术上看，恢复有两条互不相同的路线，本章两条都走：

\begin{center}
\footnotesize
\begin{tabular}{l>{\raggedright\arraybackslash}p{3.7cm}>{\raggedright\arraybackslash}p{5.4cm}l}
\toprule
路径 & 所需信息 & 核心工具 & 本章位置\\
\midrule
微分路径（参数方法） & 需求函数解析式、可微 & 路径积分、Frobenius 条件 & \S 6.2--6.5\\
非参数路径 & 有限观测点 $\{(p^t,x^t)\}$ & 显示偏好、Afriat 定理、线性规划 & \S 6.6\\
\bottomrule
\end{tabular}
\end{center}
\figexplan{两条路径的分工：微分路径假定知道整个函数，可以得到闭式的 $e$ 与 $h$，代价是需要可微性与具体的函数形式；非参数路径只用有限数据点，不需要任何函数形式，但结论是"存在一个效用函数"而非"这个效用函数"，且不能直接给出 $e$ 的解析式。两条路径在理论上相容：若数据来自一个可积的需求系统，则它必满足 GARP。}

\begin{remark}[一个必须先说清的警告：局部与全局]
可积性定理是\emph{局部}陈述：在价格空间的某个凸（或星形）邻域上，一组满足条件的偏导数可以积出一个支出函数。
把它推广到整个 $\R^n_{++}$ 需要额外的技术条件（如解曲线不逃逸、凸性在整体上成立）。
本讲按教科书传统在"相关区域"内陈述定理，并把技术边界显式标出，避免出现"定理说了全局而证明只做了局部"这类常见疏漏。
第 5 章对"相关区间"的处理（$\bar u\in(u(0),\sup u)$）出于完全相同的谨慎。
\end{remark}

\begin{remark}[术语小注]
文献中 \emph{integrability}（可积性）指"需求函数可积出支出函数"；\emph{rationalizability}（可理性化）指
"数据可由某个偏好生成"；\emph{integrability conditions} 与 \emph{Slutsky conditions} 常混用，
但严格说来前者包含后者加上可微性与边界条件。McKenzie（1957）曾指出一个方法论要点：
Slutsky 条件可以不借助任何效用概念、直接在需求函数上表述并赋予行为解释，
这正是后来显示偏好纲领的技术先声。本讲采纳这一立场——先写条件，再谈偏好。
\end{remark}
